\documentclass[10pt,a4paper]{amsart}
\usepackage[margin=19mm]{geometry}
\usepackage{amsmath,amssymb,mathtools}
\usepackage[scr=rsfs,cal=euler]{mathalfa}
\usepackage{xfrac}
\usepackage[unicode,hidelinks,bookmarks=false]{hyperref}
\newcommand{\ie}{i.e.\ }
\newcommand{\cf}{cf.\ }
\newcommand{\resp}{respectively\ }
\newcommand{\Subsection}{\subsection}
\providecommand{\nobreakdash}{\nobreak\hbox{-}}
\setlength{\emergencystretch}{2em}
\multlinegap0pt
\usepackage{amssymb}
\usepackage{mathtools}
\usepackage{enumerate}
\newtheorem{lemma}{Lemma}
\title{Classifying spaces for families of virtually abelian subgroups of surface braid groups}
\author{Ramón Flores, Juan González-Meneses, Porfirio L. León Álvarez}
\date{Source: arXiv:2604.15243v2 (CC BY 4.0)}
\begin{document}
\maketitle

\noindent\emph{Source and licence.} Classifying spaces for families of virtually abelian subgroups of surface braid groups. Version arXiv:2604.15243v2 (CC BY 4.0), distributed by arXiv under the Creative Commons Attribution 4.0 International licence, http://creativecommons.org/licenses/by/4.0/, as recorded in the arXiv licence metadata for this exact version and shown on the abstract page https://arxiv.org/abs/2604.15243v2. Full licence notice: \texttt{licenses/arxiv-2604.15243-CC-BY-4.0.txt}.\par\smallskip

\noindent\textit{Editorial scope.} Counted unit: the article's lemma property:pass:subgroup (source0.tex lines 658--662). The article's complete original proof of it is delivered with it (source0.tex lines 664--686, one \texttt{proof} environment of 23 lines). No mathematics is written, repaired or re-derived: the statement and the proof are the article's own lines, in the article's own notation, and no retained line is substituted or re-flowed.  No further source-local context is needed: every cross-reference of every retained line resolves inside the two delivered spans.  Nothing was omitted from any retained span, so the package carries no omission record.  No retained line contains a citation, so the wrapper carries no bibliography and no bibliography entry.  No retained line contains a figure environment, an image-inclusion command or drawing code, so no image is delivered and the package declares no asset dependency.  Editorial formatting only, in addition: an AMS \texttt{amsart} preamble that restates the article's own declarations and macros used by the retained lines, namely the environments \texttt{lemma} on separate counters, together with the standard packages those lines need.  The retained environments print in this document as Lemma 1 in order of appearance, whereas the article prints the same items with the numbering of its own full document.  The complete native source of the article as distributed in the arXiv e-print for this exact version is bundled unchanged as \texttt{sources/198592/source0.tex}.
\begin{lemma}\label{property:pass:subgroup}
Let $G$ be a group such that for every virtually free abelian subgroup $H \leq G$
there exists a subgroup $H' \leq G$ commensurable with $H$ satisfying
$N_G[H] = N_G(H')$. Then every subgroup of $G$ has the same property.
\end{lemma}

\begin{proof}
Let $L \leq G$ and let $H \leq L$ be a virtually free abelian subgroup. By
hypothesis, there exists a subgroup $H' \leq G$ commensurable with $H$ such that
$N_G[H] = N_G(H')$. It follows that
\[
  N_L[H] = L \cap N_G[H] = L \cap N_G(H') \subseteq N_L(H' \cap L).
\]
We claim that $H' \cap L$ is commensurable with $H$. Indeed, $H \cap (H' \cap L)
= H \cap H'$, which is commensurable with $H$ by hypothesis. It remains to show
that $H \cap H'$ has finite index in $H' \cap L$. This follows from the inequality
\[
  [H' \cap L : H \cap H'] \leq [H' : H \cap H']  < \infty,
\]
which holds since $H$ and $H'$ are commensurable. This proves the claim.

Since $H' \cap L$ is commensurable with $H$ and $H \leq L$, we have $N_L[H] =
N_L[H' \cap L]$. Combined with the inclusion established above,
\[
  N_L[H' \cap L] = N_L[H] \subseteq N_L(H' \cap L).
\]
Since the reverse inclusion $N_L(H' \cap L) \subseteq N_L[H' \cap L]$ always
holds, we conclude that $N_L(H' \cap L) = N_L[H' \cap L]$, as required.
\end{proof}

\end{document}
